% titel: Blatt 5, A2: Elektronenkonfiguration von Ionen

\subsection*{Gegeben und gesucht}
Ionen $\mathrm{Ni^{2+}, Cu^+, Fe^{3+}, Al^{3+}, I^-}$. Gesucht: vollständige und
Kurzschreibweise.

\subsection*{Lösung}
\begin{center}
\begin{tabular}{@{}lll@{}}
  Ion & Kurzschreibweise & vollständige Schreibweise \\
  \hline
  $\mathrm{Ni^{2+}}$ & $[\mathrm{Ar}]\,3d^8$ &
     $1s^2\,2s^2\,2p^6\,3s^2\,3p^6\,3d^8$ \\
  $\mathrm{Cu^+}$   & $[\mathrm{Ar}]\,3d^{10}$ &
     $1s^2\,2s^2\,2p^6\,3s^2\,3p^6\,3d^{10}$ \\
  $\mathrm{Fe^{3+}}$& $[\mathrm{Ar}]\,3d^5$ &
     $1s^2\,2s^2\,2p^6\,3s^2\,3p^6\,3d^5$ \\
  $\mathrm{Al^{3+}}$& $[\mathrm{Ne}]$ &
     $1s^2\,2s^2\,2p^6$ \\
  $\mathrm{I^-}$    & $[\mathrm{Kr}]\,4d^{10}\,5s^2\,5p^6$ &
     $1s^2\,2s^2\,2p^6\,3s^2\,3p^6\,3d^{10}\,4s^2\,4p^6\,4d^{10}\,5s^2\,5p^6$ \\
\end{tabular}
\end{center}

\subsection*{Begründung (bei Übergangsmetallen)}
Neutral: $\mathrm{Ni}\ [\mathrm{Ar}]3d^8\,4s^2$ --- erst die $4s$-Elektronen
abgeben, das $3d^8$ bleibt. $\mathrm{Cu}\ [\mathrm{Ar}]3d^{10}\,4s^1$: da die
$3d^{10}$ besonders stabil sind, gibt $\mathrm{Cu^+}$ nur das eine $4s$-Elektron
ab und behält das volle $d^{10}$. $\mathrm{Fe}^{3+}$: aus
$[\mathrm{Ar}]3d^6\,4s^2$ erst $4s^2$, dann ein $3d$-Elektron.

\begin{ergebnis}
$\mathrm{Ni^{2+}}=[\mathrm{Ar}]3d^8,\ \mathrm{Cu^+}=[\mathrm{Ar}]3d^{10},\
\mathrm{Fe^{3+}}=[\mathrm{Ar}]3d^5,\ \mathrm{Al^{3+}}=[\mathrm{Ne}],\
\mathrm{I^-}=[\mathrm{Kr}]4d^{10}5s^2\,5p^6.$
\end{ergebnis}

\begin{kontrolle}
Elektronenzahl je Ion: Ni $28-2=26$, Cu $29-1=28$, Fe $26-3=23$, Al $13-3=10$,
I $53+1=54$ --- mit der Summe der Exponenten je Zelle verglichen stimmt jede
Zeile.
\end{kontrolle}

\begin{merksatz}
Kationen: Elektronen nach \textit{steigender Hauptquantenzahl} entfernen
(erst $ns$, dann $(n-1)d$); stabile Halb-/Vollbesetzungen ($d^{10}$, $d^5$)
sind ein starkes Kriterium.
\end{merksatz}
